\subsection{赛题要求与完成情况}
系统在Ti60上完成了从摄像头原始像素到风格画面的处理。SC431HAI输出1920×1080 RAW10，FPGA将其转换为1280×720 RGB565图像。卷积网络在这一分辨率上逐像素生成风格结果，HDMI提供全屏原图、全屏风格和左原右风格三种视图。

\begin{table}[H]\centering\small
\caption{实时绘画赛题的设计要求与系统实现}\label{tab:contest}
\begin{tabularx}{\linewidth}{@{}p{.20\linewidth}p{.31\linewidth}X@{}}\toprule
功能 & 赛题要求 & 系统实现及已有结果\\\midrule
采集与预处理 & MIPI CSI-2，至少640×480、30帧/秒 & 四通道接收RAW10，完成包校验、Bayer插值及720p缩放。每发布一幅好帧更新采集计数。\\
外部帧缓存 & 双缓冲保存完整帧，保障采集连续 & DDR3设置三个原图槽、三个输出槽。采集、推理和显示分别持有所用帧，记录跳帧、坏帧和欠载。\\
卷积加速与控制 & 3×3或5×5卷积、激活和归一化，总线读写DDR，RISC-V调度 & FPGA实现3×3深度卷积、1×1通道混合和饱和激活，BN折叠入卷积。RISC-V配置摄像头、加载参数，管理作业与完成中断。\\
显示与对比 & 720p60或1080p30，原画对比及FPS叠加 & HDMI使用74.25\,MHz像素时钟输出720p60。左右两侧取自同一源帧，屏显C、I分别给出采集和推理帧率。\\
处理速度 & 基础目标为640×480稳定超过15帧/秒，高阶目标为720p达到30帧/秒 & 720p槽0对比模式的CNN/DMA整帧仿真处理921600像素，耗时25.585\,ms，折合39.085帧/秒。\\
多风格与模型规模 & 至少两种风格，切换低于100\,ms且不丢帧，模型小于500\,KB & 十组风格共用一个加速器。每风格567个可训练参数、396个卷积系数，十组部署参数共7360字节。支持按键及手势选画。\\\bottomrule
\end{tabularx}\end{table}
整机持续帧率、丢帧数和按键切换延迟尚待现场实测。
